%% chapters/03-figures-tables-code.tex
\chapter{Figures, Tables and Code}\label{chap:figures}

This chapter covers floating material (figures and tables), source code and
mathematics.

\section{Figures}\label{sec:figures}

Put image files in the \file{figures/} folder. Vector PDF files are
recommended because they stay sharp at any size and are small; PNG and JPG
also work. Thanks to \cmd{graphicspath}, only the file name is needed:

\begin{lstlisting}[language={[LaTeX]TeX}]
\begin{figure}[tb]
  \centering
  \includegraphics[width=\linewidth]{sample-plot}
  \caption{A short, descriptive caption.}
  \label{fig:sample}
\end{figure}
\end{lstlisting}

\begin{figure}[tb]
  \centering
  \includegraphics[width=\linewidth]{sample-plot}
  \caption{A sample plot, included from \file{figures/sample-plot.pdf}. Its
           \LaTeX{} source is kept in \file{figures/src/}.}
  \label{fig:sample}
\end{figure}

\Cref{fig:sample} fits the width of one column. A figure that needs more room
can span the full page width with the starred environment, as in
\cref{fig:wide}. In two-column mode such floats are placed at the top of a
page.

\begin{figure*}[tb]
  \centering
  \includegraphics[width=0.6\textwidth]{sample-plot}
  \caption{The same plot in a \texttt{figure*} environment, which spans both
           columns.}
  \label{fig:wide}
\end{figure*}

\section{Tables}\label{sec:tables}

Tables use the \pkg{booktabs} rules, which look better than full grids:
horizontal lines only, and no vertical lines. The \texttt{X} column type of
\pkg{tabularx} lets a column absorb the remaining width and wrap its text.

\begin{table}[tb]
  \centering
  \caption{Environments provided by the template.}
  \label{tab:environments}
  \begin{tabularx}{\linewidth}{@{}lX@{}}
    \toprule
    Environment & Purpose \\
    \midrule
    \texttt{note}       & Supplementary information \\
    \texttt{tip}        & Recommendations and shortcuts \\
    \texttt{warning}    & Pitfalls and common mistakes \\
    \texttt{definition} & Formal definitions, numbered per chapter \\
    \texttt{example}    & Worked examples, numbered per chapter \\
    \texttt{exercise}   & Exercises for the reader \\
    \bottomrule
  \end{tabularx}
\end{table}

\Cref{tab:environments} lists the environments introduced in
\cref{chap:writing-content}.

\section{Source Code}\label{sec:code}

Code is typeset with the \pkg{listings} package. Choose the language for each
listing, and add a caption and a label to make it citable
(\cref{lst:python}):

\begin{lstlisting}[language=Python, caption={A short Python function.},
                   label={lst:python}, numbers=left, xleftmargin=2.2em]
def average(values):
    """Return the arithmetic mean of a list of numbers."""
    if not values:
        raise ValueError("values must not be empty")
    return sum(values) / len(values)

print(average([2, 4, 6]))  # prints 4.0
\end{lstlisting}

For short fragments inside a sentence use \cmd{code}, as in
\code|average(data)|, or \cmd{texttt} for plain monospaced text. To include
a whole source file, use \cmd{lstinputlisting}.

\begin{warning}
The \cmd{code} command cannot be used inside the argument of another
command, such as \cmd{section} or \cmd{footnote}. Use \cmd{texttt} there.
\end{warning}

\section{Mathematics}\label{sec:math}

Inline mathematics is written between dollar signs, for example
$a^2 + b^2 = c^2$. Numbered equations use the \texttt{equation} environment
and can be referenced by their label, as in \cref{eq:quadratic}:

\begin{equation}\label{eq:quadratic}
  x = \frac{-b \pm \sqrt{b^{2} - 4ac}}{2a}
\end{equation}

Several aligned equations use \texttt{align}:

\begin{align}
  f(x)  &= x^2 + 1, \label{eq:f} \\
  f'(x) &= 2x.      \label{eq:fprime}
\end{align}
